\chapter{The fundamental group: generalities}
\label{V}
% original p. 105

\setcounter{section}{-1}
\section{Introduction}

The present Seminar is the continuation of the 1960 Seminar. We
refer to the latter by symbols such as I~\Ref{I.9.7},
which means: S\'eminaire de G\'eom\'etrie Alg\'ebrique,
Expos\'e~I, \No9.7. The numbers of the 1961 expos\'es will follow
those of 1960. We refer to the \'El\'ements de
G\'eom\'etrie Alg\'ebrique of Dieudonn\'e-Grothendieck by
symbols such as (EGA~I 8.7.3).

The present expos\'e summarizes (with slight
complements) the last expos\'es of 1960, which had not
been written up.

As in 1961, we shall as a general rule restrict ourselves
to locally noetherian preschemes, although this restriction is
often inessential. In
expos\'e \Ref{VI} we shall take for granted the theory of faithfully flat descent,
summarized in S\'eminaire Bourbaki \No190. If appropriate,
we shall give a more detailed exposition of it in a
later expos\'e\footnote{\Cf Exp\ptbl \Ref{VI} and Exp\ptbl \Ref{VIII}}, once
the reader has had the opportunity to convince himself of the usefulness of
this technique for the theory of the fundamental group.

\section{Prescheme with a finite group of operators,
quotient prescheme}
\label{V.1}

Let $X$ be a prescheme, $G$ a finite group operating on~$X$ by automorphisms, on the right to fix ideas. If $X$
is affine with ring $A$, $G$ therefore operates by automorphisms on
the left on~$A$.

For every prescheme $Z$, $G$ operates on the left on
the set $\Hom(X,Z)$,
one can therefore consider the set
$$
\Hom(X,Z)^G
$$
of morphisms invariant under $G$. It depends functorially on
$Z$,
one can ask whether this functor is ``representable'',
\ie isomorphic to a functor $Z\mto\Hom(Y,Z)$.
% original p. 106
This means that one can find a prescheme $Y$, and a
morphism invariant under $G$
$$
p\colon X\to Y
$$
such that for every $Z$, the corresponding map $g\mto gp$
$$
\Hom(Y,Z)\to \Hom(X,Z)^G
$$
is bijective. One then says that $(Y,p)$ is a
\emph{quotient prescheme}
of~$X$
by~$G$
(it is determined up to unique isomorphism).

\begin{proposition}
\label{V.1.1}
Let $A$ be a ring on which the finite group $G$ operates on
the left, $B=A^G$ the subring of invariants of~$A$, $X=\Spec(A)$ and
$Y=\Spec(B)$, $p\colon X\to Y$ the canonical morphism (obviously
invariant under $G$). Then
\begin{enumerate}
\item[(i)] $A$ is integral over~$B$, \ie $p$ is an
\emph{integral} morphism.
\item[(ii)] The morphism $p$ is surjective, its fibers are the
orbits of~$G$, the topology of~$Y$ is the quotient of that of~$X$.
\item[(iii)] Let $x\in X$, $y=p(x)$, $G_x$ the stabilizer of~$x$; then
$\kres(x)$ is a quasi-Galois algebraic extension of
$\kres(y)$ and the canonical map of~$G_x$ into the group
$\Gal(\kres(x)/\kres(y))$ of $\kres(y)$-automorphisms of
$\kres(x)$
is
surjective.
\item[(iv)] $(Y,p)$ is a quotient prescheme of~$X$ by $G$.
\end{enumerate}
\end{proposition}

Statements
(i), (ii), (iii)
are well known in commutative
algebra\footnote{\Cf N. Bourbaki, Alg. Comm. Chap\ptbl 5, \S 1 and \S
2, th\ptbl 2.} and are included only as a reminder, except for the assertion
about the topology, which comes from the following general fact,
an easy consequence of the Cohen-Seidenberg theorem: an
integral morphism is closed (\ie transforms closed sets into
closed sets). Let us note at once:

\begin{corollary}
\label{V.1.2}
Under the preceding conditions, the natural homomorphism
$$\cal{O}_Y{\to} p_*(\cal{O}_X)^G$$ is an isomorphism.
\end{corollary}

This follows at once from the formula
$$
(S^{-1}A)^G = S^{-1}(A^G)
$$
valid for every multiplicatively stable subset $S$ of~$B=A^G$
(a formula which extends to modules, and which can be stated more generally for
a change of base $A\to A'$ which is \emph{flat}),
applied
to the case where $S$ is generated by an element $f$ of~$B$.

Assertion (ii) and cor\ptbl \Ref{V.1.2} easily imply (iv); more
generally, one will have the following:

\begin{proposition}
\label{V.1.3}
Let
% original p. 107
$X$ be a prescheme with a finite group of automorphisms
$G$, $p\colon{}X\to{}Y$ an invariant affine morphism such that
$\cal{O}_Y\isomto p_*(\cal{O}_X)^G$. Then the conclusions
\textup{(i), (ii), (iii), (iv)}
of~\Ref{V.1.1} are still valid.
\end{proposition}

Indeed, for
(i), (ii), (iii),
one may assume $Y$, hence $X$, affine, and
if $B,A$ are their rings, the hypothesis implies $B=A^G$, it
suffices to apply~\Ref{V.1.1}. For (iv), one uses (ii) and
$\cal{O}_Y=p_*(\cal{O}_X)^G$.

\begin{corollary}
\label{V.1.4}
Under the conditions of~\Ref{V.1.3}, for every open set $U$ of~$Y$, $U$
is a quotient of~$X|U=p^{-1}(U)$ by $G$.
\end{corollary}

Indeed, $p^{-1}(U)\to{}U$ induced by $p$ satisfies the same
hypotheses as~$p$.

If now $X$ is a $Z$-prescheme and the operations
of~$G$ are $Z$-automorphisms, then by (iv) $Y$ is a
$Z$-prescheme. That said:

\begin{corollary}
\label{V.1.5}
For $X$ to be affine \resp separated over~$Z$, it is necessary and
sufficient that~$Y$ be so. If $X$ is of finite type over~$Z$, it is
\emph{finite} over~$Y$; if moreover $Z$ is locally noetherian, $Y$
is of finite type over~$Z$.
\end{corollary}

Since $X$ is affine and a fortiori separated over~$Y$, if $Y$ is
affine \resp separated over~$Z$, so
is~$X$. Conversely, suppose $X$ affine over~$Z$, let us prove that
$Y$ is: thanks to~\Ref{V.1.4} one may assume $Z$ affine, and
one is reduced to proving that if $X$ is affine, $Y$~is, which
follows from the explicit determination of~$Y$ as
$\Spec(A)^G$ made in~\Ref{V.1.1}. Likewise, since
$p\colon{}X\to{}Y$ is integral, hence universally closed, and
surjective, it follows that if $X$ is separated over~$Z$, so is
$Y$ (a lemma to be singled out!); indeed, in the diagram
$$
\xymatrix@C=1.5cm{ X \times_ZX\ar[r]^-{p\times_Zp} & Y\times_ZY \\ X
\ar[u]^{\Delta_{X/Z}}\ar[r]^{p} & Y \ar[u]_{\Delta_{Y/Z}} }
$$
the morphism $X\times_ZX\to Y\times_ZY$ is closed, hence transforms
the (closed) diagonal of~$X\times_ZX$ into a closed subset of
$Y\times_ZY$, which moreover is none other than the diagonal of the
latter since $p$ is surjective.
--- If $X$ is of finite type over~$Z$, it is a fortiori so over~$Y$, hence
it is finite over~$Y$ (since it is already integral over~$Y$). Suppose moreover $Z$ locally noetherian, let us prove that
$Y$ is of finite type over~$Z$. Thanks to~\Ref{V.1.4} one may
assume $Z$ affine. Since the topological space $X$ is quasi-compact
and $p\colon X\to Y$ is surjective, $Y$ is
likewise
quasi-compact, hence a finite union of
% original p. 108
affine open sets, and by~\Ref{V.1.4} one is reduced to the case where $Y$ is
affine, hence $X$ affine. But then the ring $A$ of~$X$ is an
algebra of finite type over the ring $C$ of~$Z$, which is
noetherian, and it is known that $B=A^G$ is then likewise
an algebra of finite type over~$C$ (for $A$ will be integral, hence
finite over a subalgebra $B'$ of~$B$ of finite type over~$C$, hence,
since $B'$ is noetherian, $B$ is likewise
finite
over~$B'$, hence of finite type over~$C$).

\begin{corollary}
\label{V.1.6}
For $X$ to be affine \resp a scheme, it is necessary and sufficient that $Y$ be
so.
\end{corollary}

\begin{definition}
\label{V.1.7}
Let
% original p. 109
$X$ be a prescheme on which a finite group $G$ operates
on the right. One says that $G$ \emph{operates admissibly}
if there exists a morphism $p\colon X\to Y$ having the properties
of~\Ref{V.1.3} (which implies that $X/G$ exists and is isomorphic
to~$Y$).
\end{definition}

\begin{proposition}
\label{V.1.8}
Let $X$ be a prescheme on which the finite group $G$ operates
on the right. For $G$ to operate admissibly, it is necessary
and sufficient that $X$ be a union of affine open sets invariant under
$G$, or again that every orbit of~$G$ in $X$ be contained
in an affine open set.
\end{proposition}

This last condition is obviously implied by the
first, and in its turn it implies it; for let $T$ be an
orbit of~$G$, $U$ an affine open set containing it; the intersection
of the transforms of~$U$ by the
$g$ in $G$
is then an open set $U'$ stable under~$G$, containing $T$ and contained in
the affine open set $U$. Since in $U$ every finite subset has a
fundamental system of affine open neighborhoods, there exists an
affine open neighborhood $V$ of~$T$ contained in $U'$. These
transforms by the
$g$ in $G$
are therefore affine and contained in $U'$, which is \emph{separated},
hence their intersection $U''$ is an affine open set which is invariant
under $G$ and contains $T$.
--- This being established, the condition considered in~\Ref{V.1.8} is
\emph{necessary}, for one will take the inverse images $X_i$
of affine open sets $Y_i$ covering~$Y$. It is sufficient, for one
can then by~\Ref{V.1.1} construct the quotients $Y_i=X_i/G$; in
each $Y_i$ the image of~$X_i\cap X_j$ is an open set $Y_{ij}$
which is identified with $X_{ij}/G$ by~\Ref{V.1.4}, in particular one
deduces from this isomorphisms $Y_{ij}\isomto Y_{ji}$ making it possible to
glue the $Y_i$ together to construct $Y$. Serre prefers to
construct directly the quotient topological space $Y$ of~$X$ by
$G$, to put on it the sheaf $p_*(\cal{O}_X)^G$ and to verify that
$Y$ becomes a prescheme and that one is then under the
conditions of~\Ref{V.1.3}.

\setcounter{subsection}{6}

\begin{corollary}
\label{cor:V.1.7}
If $G$ operating on~$X$ is admissible, the same holds for
every subgroup $H$ of~$G$ (hence $X/H$ exists).
\end{corollary}

This can also be verified directly on the
situation~\Ref{V.1.3}, by noting that one can always assume $X$
affine over some $Z$ and the $s\in G$ operating by $Z$-automorphisms
(one takes for example $Z=Y$); indeed one has:

\begin{corollary}
\label{cor:V.1.8}
Suppose $X$ affine over~$Z$, and the operations of~$G$ to be
$Z$-auto\-mor\-phisms. Then $G$ operates on~$X$
admissibly. If $X$ is defined by a quasi-coherent sheaf
$\cal{A}$ of algebras, $Y$ is defined by the sheaf
$\cal{A}^G$ of invariants of~$\cal{A}$ under~$G$.
\end{corollary}

\begin{proposition}
\label{V.1.9}
Suppose that $G$ operates admissibly on~$X$, and that
$X/G=Y$ is a prescheme over~$Z$. Consider a
change-of-base morphism $Z'\to Z$, set $X'=X\times_Z Z',
Y'=Y\times_Z Z'$, so that $G$ still operates by transport of
structure on~$X'$, the morphism $p'\colon X'\to Y'$ being
invariant. If $Z'$ is \emph{flat} over~$Z$, then $p'$ still satisfies
the hypotheses
of~\Ref{V.1.3},
\ie $\cal{O}_Y'\to
p_*'(\cal{O}_{X'})^G$ is an isomorphism ($p'$ being affine in any
case). Hence $G$ operates admissibly on~$X'$, and $(X/G)\times_Z Z'\thickapprox (X\times_Z Z')/G$.
\end{proposition}

One may obviously assume $Z=Y$, and one is reduced to the case where
moreover $Y$ and $Y'$ are affine. One must show that if $B$ is the
subring of invariants of~$G$ operating in $A$, and if $B'$ is
an algebra over~$B$ which is flat over~$B$, then $B'$ is the
subalgebra of invariants of~$A'=A\otimes_B B'$. This is
immediate, for the exact sequence
$$
0 \to B \lto{i} A \lto{j} A^{(G)}
$$
(where the last term means a power of~$A$, and where
$j(x)$ is the system of the $s\cdot x-x$, $s\in G$) remains exact under
tensoring with
$B'$.

One should take care that the flatness hypothesis was
essential for the validity of the result; in particular, if
$Y'$ is a closed subprescheme of
$Y$
(for example even a closed point of
$Y$),
$X'$ its inverse image in $X$,
then $Y'$ \emph{is not identified} in general with
$X'/G$. We shall see that it is nevertheless so if $X$ is
\'etale over~$Y$.

To finish, let us give a formalism as convenient as it is trivial. Let $Y$
be a prescheme. Since in the category of
preschemes direct sums exist, one can for every
set $E$ consider the prescheme sum of a family
$(Y_i)_{i\in E}$ of preschemes all identical to $Y$; this
prescheme will be denoted $Y\times E$.
% original p. 110
It is characterized by the formula
\begin{equation*}
\label{eq:V.1.*}
\tag{$*$} \Hom(Y\times E,Z) = \Hom(E,\Hom(Y,Z))
\end{equation*}
where the second $\Hom$ obviously denotes the set
of maps from the set $E$ into the set $\Hom(Y,Z)$. One has
a canonical morphism
$$
Y\times E\to Y
$$
making~$Y\times E$ a prescheme over~$Y$. Since
fibered products commute with direct sums (in the
category of preschemes) one will have, if $Y$ is a
prescheme over some other $Z$, for a change of base $Z'\to
Z$:
$$
(Y\times E)\times_Z Z'=(Y\times_Z Z')\times E
$$
(a formula especially useful if $Z=Y$). On the other hand, one concludes
trivially from the definition
$$
(Y\times E)\times F=Y\times(E\times F)=(Y\times E)\times_Y(Y\times F)
$$
(the last formula however resulting from the
commutativity pointed out above).

For fixed $Y$, one can regard $Y\times E$ as a functor in
$E$, with values in preschemes over~$Y$, a functor which
commutes with finite products by the preceding formula
(which makes it possible for example to associate with every ordinary group $G$ a
group scheme $Y\times G$ over~$Y$, which will be
finite over~$Y$ if $Y$ is, etc.). More generally, this
functor is ``left exact'', but we shall not need
to use this here. This functor also commutes trivially with direct
sums, and it is also ``right exact'', as one sees
at once from the defining formula \eqref{eq:V.1.*}. In
particular, if the finite group $G$ operates on the right in
the set $E$, then it operates on the right in $Y\times E$, and
one has
$$
(Y\times E)/G = Y\times (E/G)
$$
where in fact the quotient of the left-hand side satisfies the conditions
of~\Ref{V.1.3} (this is immediate).
